\documentclass[addpoints]{exam}
% \documentclass[addpoints,12pt]{exam}

\usepackage[slovene]{babel}
\usepackage{amsfonts}
\usepackage{amssymb}
\usepackage{amsmath}
\usepackage{float}
\usepackage{graphicx}
\usepackage{enumitem}
\usepackage{color}
\usepackage{eurosym}

\usepackage{pgfpages}
% \usepackage{tikz}
\usepackage{wrapfig}
\usepackage{pgfkeys}
\usepackage{pgfplots}
\usepackage{xcolor}
\usepackage{tkz-euclide}
\usepackage{xfp}




%Komentar naslednje vrstice glede na verzijo testa -- rešitve ali prostor za odgovore:
% \printanswers

% \linespread{2}


\pointsinrightmargin
\marginbonuspointname{}


\renewcommand{\solutiontitle}{\noindent\textbf{Rešitve in točkovanje:}\par\noindent}

\colorgrids
\definecolor{GridColor}{gray}{0.7}
\setlength{\gridsize}{7mm}

\extrawidth{5mm}
\setlength{\rightpointsmargin}{1.5cm}

\begin{document}

\runningfooter{}{}{}

\begin{center}
    \fbox{\fbox{\parbox{15cm}{\centering
    Peto pisno preverjanje znanja \\ 2. letnik \\ maj/junij 2026 \\ (Kvadratna funkcija; Kompleksna števila)}}}
\end{center}

\vspace{0.1in}
\makebox[\textwidth]{Ime in priimek, oddelek:\enspace\hrulefill}


\begin{center}
    \chqword{Naloga}
    \chpword{Možne točke}
    \chbpword{Dodatne točke}
    \chsword{Dosežene točke}
    \chtword{Skupaj}  
    % \combinedpointtable[h][questions]
    \combinedgradetable[h][questions]

\end{center}

\vspace{0.1in}


\begin{questions}

    % 1. vprašanje
    \question[5]
        Dani sta družina parabol $y=x^2+bx+5$ in premica $y=4x-7$.
        Kateremu pogoju mora ustrezati parameter $b$, da bo premica sekala parabolo?

            % \begin{solution}[8cm]

            % \end{solution}


    % 2. vprašanje
    \question[4]
           Nogometaš je podal žogo svojemu soigralcu.
           Na kateri višini $h$ se nahaja žoga, ko je horizontalno oddaljena za $a$ od prvega igralca, 
           lahko izračunamo po formuli $$ h(a)=-\dfrac{1}{10}a^2+2a.$$
            
            \begin{itemize}
                \item Koliko metrov proč od podajalca žoge pade žoga na tla?
                \item Koliko metrov od igralca je žoga najvišje?
                \item Na kateri višini je žoga, ko je najvišje?
            \end{itemize}

            % \begin{solution}[3cm]

            % \end{solution}


    % 3. vprašanje
    \question
        Dano je kompleksno število $$ z=\left|\left(2+i\right)^2\right|-\left(1+3i\right)\overline{1+3i}+\dfrac{25i}{3-4i}+3i^{44}. $$ 

        \begin{parts}
            
            \part[8]
            Poenostavite zapis kompleksnega števila $z$.

                % \begin{solution}[15cm]

                % \end{solution}

                
            \part[1]
            Določite realni del kompleksnega števila $z$.

                % \begin{solution}[2cm]

                % \end{solution}


            \part[1]
            Določite imaginarni del kompleksnega števila $z$.

            \part[1]
            Določite imaginarno komponento $Im(z)$ kompleksnega števila $z$.

                % \begin{solution}[2cm]

                % \end{solution}


        \end{parts}


    % \newpage

    % 4. vprašanje
    \question[6]
        V kompleksni ravnini upodobite množico točk:                     
            $$\left\{z\in\mathbb{C}; \left|z\right|\geq3 \land \left(Im(z)<Re(z)-1\right) \right\}. $$ 
        
            % \begin{solution}[12cm]

            % \end{solution}


    % 5. vprašanje
    \question[6]
        V množici kompleksnih števil rešite enačbo: $$x^3-2x^2+x=2.$$

            % \begin{solution}[8cm]

            % \end{solution}


    % dodatna naloga








\end{questions}



\end{document}